\documentclass{article}
\usepackage{graphicx}
\usepackage{listings}
\usepackage{xcolor}
\usepackage{float}
\usepackage{geometry}
\usepackage{amsmath}
\geometry{a4paper, margin=1in}

\definecolor{codegreen}{rgb}{0,0.6,0}
\definecolor{codegray}{rgb}{0.5,0.5,0.5}
\definecolor{codepurple}{rgb}{0.58,0,0.82}
\definecolor{backcolour}{rgb}{0.95,0.95,0.92}

\lstdefinestyle{mystyle}{
    backgroundcolor=\color{backcolour},   
    commentstyle=\color{codegreen},
    keywordstyle=\color{magenta},
    numberstyle=\tiny\color{codegray},
    stringstyle=\color{codepurple},
    basicstyle=\ttfamily\footnotesize,
    breakatwhitespace=false,         
    breaklines=true,                 
    captionpos=b,                    
    keepspaces=true,                 
    numbers=left,                    
    numbersep=5pt,                  
    showspaces=false,                
    showstringspaces=false,
    showtabs=false,                  
    tabsize=2
}

\lstset{style=mystyle}

\title{Compiler Design Lab Report 1\\Lexical Analysis Implementation}
\author{Abdullah Al Jubaer Gem\\ID: 210041226\\Section: 2B}
\date{\today}

\begin{document}

\maketitle

\section{Introduction}
The objective of this lab is to implement a Lexical Analyzer (Scanner) that reads a source program and identifies valid tokens such as keywords, identifiers, numbers, operators, and delimiters.

\section{Token Definitions}
The following categories of tokens are defined in the lexical analyzer:

\begin{itemize}
    \item \textbf{Keywords:} Reserved words used by the language syntax. 
    \\ \texttt{int, float, if, else, while, for, return, var}
    
    \item \textbf{Identifiers:} Names for variables, functions, etc. Defined by the regex:
    \\ \texttt{[a-zA-Z\_][a-zA-Z0-9\_]*}
    
    \item \textbf{Numbers:} 
    \begin{itemize}
        \item \textbf{Integers:} \texttt{[0-9]+}
        \item \textbf{Floats:} \texttt{[0-9]+\textbackslash.[0-9]+}
    \end{itemize}
    
    \item \textbf{Operators:} Mathematical and logical symbols.
    \\ \texttt{+, -, *, /, =, ==, !=, <, >, <=, >=}
    
    \item \textbf{Delimiters:} Punctuation marks used for structural grouping.
    \\ \texttt{(, ), \{, \}, ;, ,}
    
    \item \textbf{Whitespaces:} Characters like spaces, tabs, and newlines are used as separators but are not treated as tokens.
\end{itemize}

\section{Solution Approach}
The lexical analyzer follows a character-by-character scanning approach:
\begin{enumerate}
    \item \textbf{File Reading:} The input file is opened and read one character at a time.
    \item \textbf{Lexeme Formation:} Characters are accumulated into a buffer (lexeme) until a breakpoint (delimiter, operator, or whitespace) is encountered.
    \item \textbf{Look-ahead Logic:} For multi-character operators (e.g., \texttt{==} vs \texttt{=}), the scanner performs a look-ahead to check if the next character forms a valid combined operator.
    \item \textbf{Token Categorization:} 
    \begin{itemize}
        \item Keywords are checked against a predefined list.
        \item Identifiers and numbers are validated using Regular Expressions (Regex).
        \item Operators and delimiters are matched directly against predefined sets.
    \end{itemize}
    \item \textbf{Error Handling:} Any sequence of characters that does not match the predefined patterns is categorized as an \texttt{INVALID} token.
\end{enumerate}

\section{Implementation}
The implementation is done in Python. The source code is provided below:

\begin{lstlisting}[language=Python, caption=Lexical Analyzer Source Code]
import re

keywords = ["int", "float", "if", "else", "while", "for", "return", "var"]
whitespace = [" ", "\n", "\t"]
delimiters = ["(", ")", "{", "}", ";", ","]
operators = ["+", "-", "*", "/", "=", "==", "!=", "<", ">", "<=", ">="]

identifier_pattern = r"[a-zA-Z_][a-zA-Z0-9_]*"
integer_pattern = r"[0-9]+"
float_pattern = r"[0-9]+\.[0-9]+"

def is_keyword(token):
    return token in keywords

def is_delimiter(token):
    return token in delimiters

def is_operator(token):
    return token in operators

def is_identifier(token):
    return re.fullmatch(identifier_pattern, token) is not None

def is_integer(token):
    return re.fullmatch(integer_pattern, token) is not None

def is_float(token):
    return re.fullmatch(float_pattern, token) is not None

def print_token(token_type, lexeme):
    print(f"<{token_type}, {lexeme}>")

def identify_token(lexeme):
    if is_keyword(lexeme):
        print_token("KEYWORD", lexeme)
    elif is_float(lexeme):
        print_token("NUMBER", lexeme)
    elif is_integer(lexeme):
        print_token("NUMBER", lexeme)
    elif is_identifier(lexeme):
        print_token("IDENTIFIER", lexeme)
    else:
        print_token("INVALID", lexeme)

def lexical_analyzer(filename):
    lexeme = ""
    pending_char = ""

    with open(filename, "r") as file:
        while True:
            if pending_char != "":
                char = pending_char
                pending_char = ""
            else:
                char = file.read(1)

            if char == "":
                if lexeme != "":
                    identify_token(lexeme)
                break

            if char in whitespace:
                if lexeme != "":
                    identify_token(lexeme)
                    lexeme = ""
                continue

            if char in delimiters:
                if lexeme != "":
                    identify_token(lexeme)
                    lexeme = ""
                print_token("DELIMITER", char)
                continue

            if char in operators:
                if lexeme != "":
                    identify_token(lexeme)
                    lexeme = ""

                next_char = file.read(1)
                possible_operator = char + next_char

                if possible_operator in operators:
                    print_token("OPERATOR", possible_operator)
                else:
                    if char in operators:
                        print_token("OPERATOR", char)
                    else:
                        print_token("INVALID", char)
                    pending_char = next_char
                continue

            if char.isalnum() or char == "_" or char == ".":
                lexeme += char
                continue

            if lexeme != "":
                identify_token(lexeme)
                lexeme = ""

            print_token("INVALID", char)

if __name__ == "__main__":
    lexical_analyzer("input.txt")
\end{lstlisting}

\section{Sample Input and Output}

\subsection{Input (\texttt{input.txt})}
\begin{lstlisting}
int a = 10
float b = a + 10
var x = (b / 3)
\end{lstlisting}

\subsection{Output Log}
\begin{lstlisting}
<KEYWORD, int>
<IDENTIFIER, a>
<OPERATOR, =>
<NUMBER, 10>
<KEYWORD, float>
<IDENTIFIER, b>
<OPERATOR, =>
<IDENTIFIER, a>
<OPERATOR, +>
<NUMBER, 10>
<KEYWORD, var>
<IDENTIFIER, x>
<OPERATOR, =>
<DELIMITER, (>
<IDENTIFIER, b>
<OPERATOR, />
<NUMBER, 3>
<DELIMITER, )>
\end{lstlisting}

\section{Conclusion}
The lexical analyzer successfully correctly identifies and tokenizes the given input program according to the specified rules. It handles keywords, identifiers, numbers (both integer and float), operators (including multi-character ones), and delimiters while effectively ignoring whitespaces.

\end{document}
